% SYN-ACK WIRE -- Rule Book
% Build: make            (regenerates ../style/theme.tex from theme/theme.toml, then runs pdflatex twice)
% The rules text is shared with the platform manuals: ../rules/*.tex.
%
% All names, vocabulary and rule numbers come from theme/theme.toml via theme.tex:
%   \Game{key}  [game] value          \Rule{key}  [rules] number (formatted)
%   \V{key}     term, singular        \Vp{key}    plural
%   \VC{key}    Capitalized           \VCp{key}   Capitalized plural
%   \Syn{i}{field}  syndicate i (0-6): name short letter coin tier alias flavor; color "syn<i>"
%   \Glyph{key} board glyph (empty, spike)
\documentclass[11pt,letterpaper]{article}

\usepackage[T1]{fontenc}
\usepackage[sfdefault,book,lining,tabular]{FiraSans}
\usepackage[scale=0.88]{FiraMono}
\usepackage{textcomp}
\usepackage[margin=0.85in,top=0.9in,bottom=0.9in]{geometry}
\usepackage{xcolor}
\usepackage{tikz}
\usepackage{booktabs}
\usepackage{tabularx}
\usepackage{array}
\usepackage{colortbl}
\usepackage{enumitem}
\usepackage{multicol}
\usepackage{needspace}
\usepackage{titlesec}
\usepackage{fancyhdr}
\usepackage[most]{tcolorbox}
\usepackage[hidelinks]{hyperref}

% ------------------------------------------------------------------
% Theme lookup
% ------------------------------------------------------------------
\makeatletter
\newcommand{\th@get}[1]{\ifcsname #1\endcsname\csname #1\endcsname\else
  \PackageError{theme}{Unknown theme key '#1'}{Check theme/theme.toml}\fi}
\newcommand{\Game}[1]{\th@get{th@game@#1}}
\newcommand{\Rule}[1]{\th@get{th@rule@#1}}
\newcommand{\Glyph}[1]{\th@get{th@glyph@#1}}
\newcommand{\V}[1]{\th@get{th@v@#1}}
\newcommand{\Vp}[1]{\th@get{th@vp@#1}}
\newcommand{\VC}[1]{\th@get{th@vc@#1}}
\newcommand{\VCp}[1]{\th@get{th@vcp@#1}}
\newcommand{\Syn}[2]{\th@get{th@syn@#1@#2}}
\newcommand{\TierCoins}[1]{\th@get{th@tiercoins@#1}}
\makeatother
\input{theme.tex}

% ------------------------------------------------------------------
% Colors
% ------------------------------------------------------------------
\definecolor{ink}{HTML}{1B1F2A}
\definecolor{neon}{HTML}{00A6C8}
\definecolor{hot}{HTML}{D6246E}
\definecolor{paper}{HTML}{F3F6F9}
\definecolor{rule}{HTML}{B8C2CC}
\definecolor{spike}{HTML}{9AA5B1}

\color{ink}

% ------------------------------------------------------------------
% Macros
% ------------------------------------------------------------------
\newcommand{\Y}[1]{\Game{currency}#1}
\newcommand{\node}[1]{\texttt{#1}}
\newcommand{\term}[1]{\textbf{\color{hot}#1}}
\newcommand{\game}{\Game{title}}
\newcommand{\atscreen}[1]{}                 % where a platform manual shows its screen (docs/manual)

\setlist{itemsep=2pt,topsep=4pt}

\titleformat{\section}{\Needspace{8\baselineskip}\Large\bfseries\color{neon}}{\texttt{\thesection}}{0.6em}{}[{\color{rule}\titlerule}]
\titleformat{\subsection}{\large\bfseries\color{ink}}{}{0em}{}
\titlespacing*{\section}{0pt}{16pt}{8pt}
\titlespacing*{\subsection}{0pt}{12pt}{4pt}

\pagestyle{fancy}
\fancyhf{}
\renewcommand{\headrulewidth}{0pt}
\fancyfoot[L]{\small\color{spike}\texttt{\Game{title} // RULE BOOK}}
\fancyfoot[R]{\small\color{spike}\texttt{[\thepage]}}

\newtcolorbox{example}[1][]{enhanced,breakable,colback=paper,colframe=neon,
  boxrule=0.6pt,arc=2pt,left=8pt,right=8pt,top=6pt,bottom=6pt,
  fonttitle=\bfseries\small,coltitle=white,title={EXAMPLE#1}}
\newtcolorbox{note}[1]{enhanced,breakable,colback=white,colframe=hot,
  boxrule=0.6pt,arc=2pt,left=8pt,right=8pt,top=6pt,bottom=6pt,
  fonttitle=\bfseries\small,coltitle=white,title={#1}}
\newtcolorbox{terminal}{enhanced,colback=ink,colframe=ink,arc=2pt,
  left=8pt,right=8pt,top=6pt,bottom=6pt,fontupper=\ttfamily\small\color{neon}}

% ------------------------------------------------------------------
% Grid drawing (TikZ). Columns A-L left to right, rows 0-8 top to bottom.
% The board shows each syndicate as its letter; color is a bonus where the platform allows.
%   \hive{col}{row}{syndicate id}   \spk{col}{row}   \target{col}{row}{marker}
% ------------------------------------------------------------------
\newcommand{\gridframe}{%
  \foreach \c [count=\i from 0] in {A,B,C,D,E,F,G,H,I,J,K,L}
    \node[font=\ttfamily\small,color=spike] at (\i+0.5,0.45) {\c};
  \foreach \r in {0,...,8}
    \node[font=\ttfamily\small,color=spike] at (-0.45,-\r-0.5) {\r};
  \foreach \c in {0,...,11} \foreach \r in {0,...,8}
    \draw[rule,fill=white] (\c+0.06,-\r-0.06) rectangle (\c+0.94,-\r-0.94);
}
\newcommand{\hive}[3]{%
  \draw[ink,line width=0.4pt,fill=syn#3] (#1+0.06,-#2-0.06) rectangle (#1+0.94,-#2-0.94);
  \node[font=\ttfamily\small\bfseries,color=white] at (#1+0.5,-#2-0.5) {\Syn{#3}{letter}};
}
\newcommand{\spk}[2]{%
  \draw[ink,line width=0.4pt,fill=spike] (#1+0.06,-#2-0.06) rectangle (#1+0.94,-#2-0.94);
  \node[font=\ttfamily\small\bfseries,color=white] at (#1+0.5,-#2-0.5) {\Glyph{spike}};
}
\newcommand{\target}[3]{%
  \draw[hot,line width=1.2pt,dashed] (#1+0.06,-#2-0.06) rectangle (#1+0.94,-#2-0.94);
  \node[font=\ttfamily\small\bfseries,color=hot] at (#1+0.5,-#2-0.5) {#3};
}

% ==================================================================
\begin{document}

% ------------------------------------------------------------------
% Cover
% ------------------------------------------------------------------
\begin{titlepage}
\pagecolor{ink}\color{white}
\centering
\vspace*{1.2in}
{\fontsize{46}{50}\selectfont\bfseries\ttfamily\color{neon} \Game{title}\par}
\vspace{0.25in}
{\Large\color{hot}\ttfamily [ RULE BOOK ]\par}
\vspace{0.5in}
\begin{tikzpicture}[scale=0.42]
  \foreach \c in {0,...,11} \foreach \r in {0,...,8}
    \draw[spike!40!ink,fill=ink] (\c+0.08,-\r-0.08) rectangle (\c+0.92,-\r-0.92);
  \foreach \c/\r in {2/1,3/1,3/2,4/2} \fill[syn5] (\c+0.08,-\r-0.08) rectangle (\c+0.92,-\r-0.92);
  \foreach \c/\r in {7/3,8/3,8/4,9/4,9/5} \fill[syn4] (\c+0.08,-\r-0.08) rectangle (\c+0.92,-\r-0.92);
  \foreach \c/\r in {1/6,2/6,2/7} \fill[syn2] (\c+0.08,-\r-0.08) rectangle (\c+0.92,-\r-0.92);
  \foreach \c/\r in {5/5,10/1,6/7} \fill[spike] (\c+0.08,-\r-0.08) rectangle (\c+0.92,-\r-0.92);
  \draw[hot,line width=1.4pt] (5+0.08,-2-0.08) rectangle (5+0.92,-2-0.92);
\end{tikzpicture}
\vspace{0.5in}

{\large Seven syndicates. One region. Every \V{node} for sale.\par}
\vspace{0.15in}
{\color{spike} \Rule{minplayers}--\Rule{maxplayers} \VCp{playertitle} \quad//\quad networked play with bots\par}
\vfill
{\small\color{spike} An unofficial cyberpunk re-skin of \emph{Acquire} by Sid Sackson.\par}
\end{titlepage}
\nopagecolor\color{ink}

% ------------------------------------------------------------------
\section{The Network Grid}
% ------------------------------------------------------------------
\input{rules/opening.tex}

% ------------------------------------------------------------------
\section{What's in the Game}
% ------------------------------------------------------------------
\input{rules/contents.tex}

% ------------------------------------------------------------------
\section{The Seven Syndicates}
% ------------------------------------------------------------------
\input{rules/syndicates.tex}

% ------------------------------------------------------------------
\section{The Table}
% ------------------------------------------------------------------
\begin{enumerate}
  \item The first player to join the server is assigned as the host.
  \item The host can decide how many seats will be at the table (2-6).
  \item Other players
\end{enumerate}

% ------------------------------------------------------------------
\section{Setup}
% ------------------------------------------------------------------
\input{rules/setup.tex}

% ------------------------------------------------------------------
\section{Playing a Turn}
% ------------------------------------------------------------------
\input{rules/turn.tex}

% ------------------------------------------------------------------
\section{Ending the Game}
% ------------------------------------------------------------------
\input{rules/ending.tex}

% ------------------------------------------------------------------
\section{Network Play}
% ------------------------------------------------------------------
\game\ is played over the network, with each \V{player} at their own machine.
\begin{itemize}
  \item \textbf{Bots fill empty seats.} Start a game with as many humans as you have; bots take the rest.
  \item \textbf{Shot clock.} Every decision has a time limit. If yours runs out, your bot makes that one decision for you.
  \item \textbf{Dropped connection.} If you drop off the network, your bot plays your seat until you reconnect, then hands
    it back. Your \V{kit}, \Vp{coin} and \Game{currency} are waiting for you.
\end{itemize}

% ------------------------------------------------------------------
\section{Frequently Asked Questions}
% ------------------------------------------------------------------
\input{rules/faq.tex}

% ------------------------------------------------------------------
\section{Where the Money Comes From}
% ------------------------------------------------------------------
\input{rules/money.tex}

% ------------------------------------------------------------------
\section{Tips}
% ------------------------------------------------------------------
\input{rules/tips.tex}

% ------------------------------------------------------------------
\section{Glossary}
% ------------------------------------------------------------------
\input{rules/glossary.tex}

% ------------------------------------------------------------------
\clearpage
\section{Reference Chart}
% ------------------------------------------------------------------
Find the \V{hive}'s size in its tier's column, then read across for the \V{coin} price and the payouts.

\medskip
\begin{center}
\renewcommand{\arraystretch}{1.35}
\large
\begin{tabular}{ccc|r|rr}
\toprule
\multicolumn{3}{c|}{\textbf{\VC{hive} size (\Vp{node})}} & \textbf{\VC{coin}} & \multicolumn{2}{c}{\textbf{Payout}} \\
\textbf{Low} & \textbf{Mid} & \textbf{High} & \textbf{price} & \textbf{\V{root}} & \textbf{\V{sudo}} \\
\TierCoins{low} & \TierCoins{mid} & \TierCoins{high} & & & \\
\midrule
2     & --    & --    & \Y{200}   & \Y{2,000}  & \Y{1,000} \\
3     & 2     & --    & \Y{300}   & \Y{3,000}  & \Y{1,500} \\
4     & 3     & 2     & \Y{400}   & \Y{4,000}  & \Y{2,000} \\
5     & 4     & 3     & \Y{500}   & \Y{5,000}  & \Y{2,500} \\
6--10 & 5     & 4     & \Y{600}   & \Y{6,000}  & \Y{3,000} \\
11--20 & 6--10 & 5    & \Y{700}   & \Y{7,000}  & \Y{3,500} \\
21--30 & 11--20 & 6--10 & \Y{800} & \Y{8,000}  & \Y{4,000} \\
31--40 & 21--30 & 11--20 & \Y{900} & \Y{9,000} & \Y{4,500} \\
41+   & 31--40 & 21--30 & \Y{1,000} & \Y{10,000} & \Y{5,000} \\
--    & 41+   & 31--40 & \Y{1,100} & \Y{11,000} & \Y{5,500} \\
--    & --    & 41+   & \Y{1,200} & \Y{12,000} & \Y{6,000} \\
\bottomrule
\end{tabular}
\end{center}

\medskip
\begin{center}
\begin{tabular}{@{}ll@{}}
\textbf{\Rule{safesize}+ \Vp{node}} & \VC{safe}: the \V{hive} can't be \V{absorbed}. \\
\textbf{\Rule{endsize}+ \Vp{node}} & \VC{totalcompromise}: the \V{nullroute} may be called. \\
\textbf{\V{root} bonus} & 10 $\times$ \V{coin} price. \quad \textbf{\V{sudo} bonus}: 5 $\times$ \V{coin} price. \\
\end{tabular}
\end{center}

\vfill
\begin{center}
\small\color{spike}
\game\ is an unofficial fan re-skin of \emph{Acquire}, designed by Sid Sackson (1964).\\
\emph{Acquire} is a trademark of Hasbro, Inc. This project is not affiliated with or endorsed by Hasbro.
\end{center}

\end{document}
